% =====================================================================
\chapter{Demostración directa y por contrarrecíproco}
\label{tema:3}
% =====================================================================

La mayoría de los teoremas tienen la forma $\forall x\,(P(x)\Rightarrow Q(x))$. Este tema presenta las dos técnicas básicas para demostrarlos: la demostración directa, que parte de $P(x)$ y llega a $Q(x)$, y la demostración por contrarrecíproco, que parte de $\neg Q(x)$ y llega a $\neg P(x)$. Como banco de pruebas se usa la aritmética de los enteros (paridad, divisibilidad y congruencias), igual que en los capítulos~4 y~5 de~\cite{hammack}, la §3.2 de~\cite{velleman}, las §1.5–1.6 de~\cite{mcs} y el primer \emph{problem set} de CS103~\cite{cs103}.

\section{Definiciones de partida}

Una demostración solo puede usar lo que está definido o demostrado antes. Se admiten como conocidas las propiedades algebraicas habituales de $\Z$ (la suma, la resta y el producto de enteros son enteros; asociatividad, conmutatividad, distributividad) y de $\R$ (incluidas las reglas de las desigualdades), y se fijan las definiciones siguientes~\cite[§4.2]{hammack}.

\begin{definicion}[Paridad]\label{def:paridad}
Un entero $n$ es \emph{par} si existe $k\in\Z$ con $n=2k$, e \emph{impar} si existe $k\in\Z$ con $n=2k+1$. Dos enteros tienen la \emph{misma paridad} si ambos son pares o ambos impares.
\end{definicion}

\begin{definicion}[Divisibilidad]\label{def:divide}
Sean $a,b\in\Z$. Se dice que $a$ \emph{divide} a $b$, y se escribe $a\divides b$, si existe $c\in\Z$ con $b=ac$. Se dice también que $b$ es \emph{múltiplo} de $a$ y que $a$ es \emph{divisor} de $b$. Un entero $p\ge 2$ es \emph{primo} si sus únicos divisores positivos son $1$ y $p$; un entero $n\ge 2$ que no es primo es \emph{compuesto}.
\end{definicion}

\begin{teorema}[División euclídea]\label{teo:division}
Sean $a\in\Z$ y $b\in\N$. Existen enteros únicos $q$ y $r$ con $a=qb+r$ y $0\le r<b$.
\end{teorema}

La existencia se demuestra en el \cref{tema:5} (\cref{prop:division-existencia}) y la unicidad en el \cref{tema:6} (\cref{prop:division-unicidad}). Con $b=2$, el teorema dice que todo entero es par o impar, y no ambas cosas; este hecho se usará libremente.

\section{Demostración directa}

\begin{definicion}[Demostración directa]\label{def:directa}
Para demostrar $\forall x\,(P(x)\Rightarrow Q(x))$ de forma \emph{directa}, se toma un $x$ arbitrario, se supone $P(x)$ y, mediante definiciones, resultados previos y reglas de inferencia, se deduce $Q(x)$.
\end{definicion}

El esquema es siempre el mismo: «Sea $x$ [del universo] tal que $P(x)$. [\dots] Por tanto, $Q(x)$.» Antes de empezar conviene escribir qué significa la tesis según las definiciones, para saber a qué expresión hay que llegar~\cite[§4.3]{hammack}.

\begin{proposicion}\label{prop:impar-cuadrado}
Si $n$ es un entero impar, entonces $n^{2}$ es impar.
\end{proposicion}

Obsérvese que la demostración introduce un nombre nuevo, $m$, para el entero que da la forma $2m+1$: el objetivo era encontrar \emph{algún} entero con esa propiedad, y la definición de impar no exige que sea el mismo $k$ de la hipótesis.

\begin{proposicion}[Propiedades de la divisibilidad]\label{prop:divisibilidad}
Sean $a,b,c,x,y\in\Z$.
\begin{enumerate}
  \item Si $a\divides b$ y $b\divides c$, entonces $a\divides c$.
  \item Si $a\divides b$ y $a\divides c$, entonces $a\divides(bx+cy)$.
\end{enumerate}
\end{proposicion}

\addtocounter{definicion}{1}%

\section{Demostración por contrarrecíproco}

\begin{definicion}[Demostración por contrarrecíproco]\label{def:contrarreciproco}
Para demostrar $\forall x\,(P(x)\Rightarrow Q(x))$ por \emph{contrarrecíproco}, se demuestra de forma directa $\forall x\,(\neg Q(x)\Rightarrow\neg P(x))$. Es legítimo porque ambos enunciados son lógicamente equivalentes (\cref{prop:leyes}(v)).
\end{definicion}

La técnica es útil cuando la negación de la tesis aporta información más manejable que la hipótesis. Un indicio típico es que la tesis afirme algo sobre un número ($n$ es par) y la hipótesis, algo sobre una expresión de ese número ($n^{2}$ es par)~\cite[§5.1]{hammack}.

\begin{proposicion}\label{prop:cuadrado-par}
Sea $n\in\Z$. Si $n^{2}$ es par, entonces $n$ es par.
\end{proposicion}

Un intento directo habría partido de $n^{2}=2k$ y necesitaría extraer una raíz cuadrada, que no conduce a ninguna parte. El contrarrecíproco convierte el problema en el cálculo ya hecho.

\addtocounter{definicion}{1}%

\section{Congruencias}

\begin{definicion}[Congruencia]\label{def:congruencia}
Sea $n\in\N$. Dos enteros $a$ y $b$ son \emph{congruentes módulo $n$}, y se escribe $a\equiv b\pmod n$, si $n\divides(a-b)$.
\end{definicion}

Por la división euclídea, todo entero es congruente módulo $n$ con exactamente uno de $0,1,\dots,n-1$, su resto. Por ejemplo, $n$ es par si y solo si $n\equiv 0\pmod 2$.

\begin{proposicion}\label{prop:congruencias}
Sea $n\in\N$ y sean $a,b,c,d\in\Z$ con $a\equiv b\pmod n$ y $c\equiv d\pmod n$. Entonces
\begin{enumerate}
  \item $a+c\equiv b+d\pmod n$;
  \item $ac\equiv bd\pmod n$;
  \item $a^{k}\equiv b^{k}\pmod n$ para todo $k\in\N$.
\end{enumerate}
\end{proposicion}

\addtocounter{definicion}{1}%

\section{Bicondicionales y cadenas de equivalencias}

\begin{definicion}[Demostración de un bicondicional]
Para demostrar $P\Leftrightarrow Q$ se demuestran por separado $P\Rightarrow Q$ y $Q\Rightarrow P$, cada una con la técnica más conveniente (\cref{prop:leyes}(vii)). Para demostrar que varias afirmaciones $P_1,\dots,P_k$ son equivalentes basta una cadena cíclica $P_1\Rightarrow P_2\Rightarrow\dots\Rightarrow P_k\Rightarrow P_1$~\cite[§7.2]{hammack}.
\end{definicion}

\begin{proposicion}\label{prop:par-equiv}
Para $n\in\Z$, son equivalentes: (a) $n$ es par; (b) $n^{2}$ es par; (c) $n^{2}+2n$ es múltiplo de~4.
\end{proposicion}

\addtocounter{definicion}{1}%

\addtocounter{definicion}{1}%

